\cotefile{29}
\scanpage{20}
Dans ce cas, le foncteur ``réalisation géométrique'' est bien défini et induit sur chaque $\mathrm{Rev}_{\underline{R}_i}(S)$ le foncteur réalisation géométrique déjà explicité dans \hspace{3em} l'exemple précédent.

de $Z_j$, on a $Z'_j\quad Z'_i/N_{ij}$, où $Z'_i$ est l'image inverse de $Z'_j$ sur $Z_i$). Notons que les critères signalés dans l'exemple 1 pour que le foncteur $\underline{r}:\mathrm{Rev}_{\underline{R}}(S)\qquad(\mathrm{Sch})_{/S}$ soit fidèle resp. pleinement fidèle, s'étendent de façon immédiate au cas envisagé ici.

\uline{Exemple} 3 Supposons que $\underline{R}=(\underline{C},\underline{s})$ satis\unsure{\deleted{g}f|\deleted{f}g}\unsure{\deleted{q}a|\deleted{a}q}sse aux conditions suivantes: pour tout $S'\quad\mathrm{Ob}\,\underline{\mathrm{Et}}_S$, et $M\quad\mathrm{Ob}\,\underline{C}_{S'}$, le groupe des automorphismes $G$ de $M$ est fini, et pour toute flèche $M'\quad M$ de $\underline{C}_{S'}$, il existe un unique morphisme de groupe $G'=\mathrm{Aut}(M')$ dans $G=\mathrm{Aut}(M)$ tel que $M'\quad M$ soit compatible avec les opérations de $G'$, de plus si $H$ est son noyau, le morphisme $\underline{s}(M')\quad\underline{\unsure{s|\deleted{e}s|\deleted{s}e}}(M)$ induit un isomorphisme $\underline{s}(M')/H\quad\underline{s}(M)$. \deleted{Enfin} On suppose \deleted{la catégorie $\underline{C}_{S'}$ filtrante (i.e. de pour deux objets $M$, $M'$ de ladite, il existe un objet $M''$ et des flèches $M\quad M''$, $M'\quad M''$),}\inserted{\deleted{que} de plus que pour deux morphismes de $M'$ dans $M^{\deleted{\prime}}$, \deleted{les morphismes correspondants de $G'$ dans $G$ sont \unsure{conjugués,|conjugués} déduit l'un de l'autre par composition avec un}} l'un est déduit de l'autre par composition avec un élément de $G$. Enfin, on suppose qu'étant donnés deux objets $M$; $M'$ de $\underline{C}_{S'}$, il existe une famille surjective de morphismes étales $S'_i\quad S'$ telle que pour tout $i$, $M_{S'}$ et $M'_{S_i}$ soient majorés par un objet $M''_i$ de $\underline{C}_{S'_i}$. Utilisant la théorie de la descente, on voit sous ces conditions que \unsure{\deleted{d}l|\deleted{l}d}\unsure{\deleted{e}a|\deleted{a}e} foncteur réalisation géométrique \deleted{de} $\mathrm{Rev}_{\underline{R}}(S)\qquad(\mathrm{Sch})_{/S}$ est défini, et à valeurs dans les \deleted{schémas fini\unsure{s} sur} revêtements de $S$ (en supposant que les $\underline{s}(M)$ sont des revêtements, bien sûr; pour pouvoir appliquer la descente, il suffirait même que les $\underline{s}(M)$ soient affines sur $S'$). Ici encore, on a des critères pour que ce foncteur soit fidèle resp. pleinement fidèle: il suffit que pour tout $S'\quad\mathrm{Ob}\,\underline{\mathrm{Et}}_S$ et tout $M\quad\mathrm{Ob}\,\underline{C}_{S'}$, la réalisation géométrique $\underline{s}(M)$, munie du groupe d'automorphismes $G=\mathrm{Aut}(M)$, satisfasse aux conditions \inserted{correspondantes} envisagées dans l'exemple 1 pour $(Z,G)$.

\uline{Cas particuliers du type kummérien}: a) Supposons données des sections
% archive sheet number: 19
